\documentclass[11pt,a4paper]{article}
\usepackage[utf8]{inputenc}
\usepackage[T1]{fontenc}
\usepackage{amsmath,amssymb}
\usepackage{bm}
\usepackage{geometry}
\usepackage{hyperref}
\usepackage{amsthm}
\usepackage{booktabs}
\usepackage{microtype}
\newtheorem{theorem}{Theorem}
\newtheorem{definition}{Definition}
\newtheorem{remark}{Remark}
\geometry{margin=2.7cm}
\hypersetup{hidelinks}

\title{Moment--Centred Decomposition (MVDC):\\
Central Bernoulli Numbers and High--Precision Euler--Product Corrections}
\author{Ing.~Robert~Polak\\\small e--mail: robopol@gmail.com}
\date{August 11, 2026}

\begin{document}
\maketitle

\begin{abstract}
The Mean--Value Decomposition by Centre (MVDC) separates a finite product
into a chosen centre term and an exact residual.  For smooth factor
sequences, midpoint Taylor expansion of the logarithm is governed by
centred powers of the index.  These finite coefficients form a natural
\(C\)--Bernoulli family and admit exact representations through central
Bernoulli polynomials and a compact generating function.  A second,
distinct moment expansion follows from the Taylor series of
\(\log(1+x)\) when the factors, rather than their logarithms, are centred
multiplicatively.  Keeping these two constructions separate removes an
ambiguity in an earlier draft while preserving the central Bernoulli
structure of MVDC.  The factorial application gives the usual exact
Bernoulli coefficients around an analytic MVDC centre.  For Euler
products the finite prime block is evaluated exactly by its logarithmic
mean, and only the remaining infinite tail must be modelled or bounded.
\end{abstract}

\section{Exact decomposition and the role of the centre}

Let
\[
  P_n=\prod_{i=1}^{n}a_i,\qquad a_i>0,\qquad
  \ell_i:=\log a_i .
\]
For every real centre \(c\), with \(k=e^c\), define
\[
  H_c:=k^n=e^{nc},\qquad
  R_c(n):=\sum_{i=1}^{n}(\ell_i-c).
\]
Then the fundamental MVDC identity is
\begin{equation}\label{eq:exact-centre}
  \boxed{\;P_n=H_c\,e^{R_c(n)},\qquad
  \log P_n=nc+R_c(n).\;}
\end{equation}
No asymptotic assumption is used in \eqref{eq:exact-centre}.  The centre
may be selected analytically so that \(H_c\) absorbs known dominant
terms, or numerically for stability.

There is one important limiting case.  If
\[
  c=\mu_{\log}:=\frac1n\sum_{i=1}^{n}\ell_i,
\]
then \(R_{\mu_{\log}}(n)=0\) and \(H_{\mu_{\log}}=P_n\) identically.
This is a valid exact log-domain evaluation, but it leaves no nonzero
residual to approximate.  A nontrivial asymptotic application therefore
uses a centre specified independently of the exact finite product, for
example by an integral, an endpoint calculation, or known leading
asymptotics.

\section{The correct factor-moment expansion}

Higher moments enter through the factors themselves.  Write
\[
  a_i=k(1+x_i),\qquad x_i=\frac{a_i}{k}-1,
  \qquad X_r(k):=\sum_{i=1}^{n}x_i^r .
\]
If \(\rho:=\max_i|x_i|<1\), summing the Taylor series for
\(\log(1+x_i)\) gives
\begin{equation}\label{eq:factor-moments}
  \boxed{\;
  \log\frac{P_n}{k^n}
  =\sum_{r=1}^{\infty}\frac{(-1)^{r-1}}{r}\,X_r(k).\;}
\end{equation}
After truncation at \(r=M\), the remainder obeys the explicit bound
\[
  |\mathcal R_M|
  \le
  \sum_{i=1}^{n}
  \frac{|x_i|^{M+1}}{(M+1)(1-|x_i|)}.
\]
For the arithmetic-mean centre
\(k=\bar a=n^{-1}\sum_i a_i\), one has \(X_1(k)=0\), so the series starts
at \(r=2\).  This is the precise setting in which the cancellation of
the first moment follows from \(\log(1+x)\).  Logarithmic central
moments \(\sum(\ell_i-\mu_{\log})^r\) are useful descriptors of the
data, but they are not interchangeable with the factor moments
\(X_r(k)\) in \eqref{eq:factor-moments}.

\section{Central Bernoulli and Nörlund polynomials}

Ordinary Bernoulli polynomials are defined by
\[
  \frac{t e^{xt}}{e^t-1}
  =\sum_{r=0}^{\infty}B_r(x)\frac{t^r}{r!}.
\]
The generalized Bernoulli polynomials of order \(\alpha\), often
associated with Nörlund, are defined by
\begin{equation}\label{eq:norlund}
  \left(\frac{t}{e^t-1}\right)^\alpha e^{xt}
  =\sum_{r=0}^{\infty}B_r^{(\alpha)}(x)\frac{t^r}{r!}.
\end{equation}
One standard central Bernoulli family is
\begin{equation}\label{eq:central-bernoulli}
  B_r^c(x):=2^rB_r\!\left(\frac{x+1}{2}\right),
  \qquad
  \frac{t}{\sinh t}e^{xt}
  =\sum_{r=0}^{\infty}B_r^c(x)\frac{t^r}{r!}.
\end{equation}
It satisfies \(B_r^c(-x)=(-1)^rB_r^c(x)\).  More generally,
\[
  \left(\frac{t}{\sinh t}\right)^\alpha e^{xt}
  =\sum_{r=0}^{\infty}B_{r,\alpha}^c(x)\frac{t^r}{r!},
  \qquad
  B_{r,\alpha}^c(x)
  =2^rB_r^{(\alpha)}\!\left(\frac{x+\alpha}{2}\right).
\]
Thus central Bernoulli objects and their Nörlund connection are
standard; the remaining question is how they enter a particular MVDC
expansion.

\section{Finite C--Bernoulli coefficients}

Put \(m_n=(n+1)/2\).  The finite centred-index coefficients used in this
note are
\begin{definition}[Finite \(C\)--Bernoulli coefficients]
\begin{equation}\label{eq:C-definition}
  \boxed{\;
  C_r(n):=\frac1{r!}\sum_{i=1}^{n}(i-m_n)^r,\qquad r\ge0.\;}
\end{equation}
\end{definition}
This definition is purely algebraic and depends only on the geometry of
the \(n\) equally spaced indices.  Faulhaber's identity gives
\begin{equation}\label{eq:C-bernoulli}
  C_r(n)
  =
  \frac{
    B_{r+1}\!\left(\frac{n+1}{2}\right)
    -B_{r+1}\!\left(\frac{1-n}{2}\right)}
  {(r+1)r!}
  =
  \frac{B_{r+1}^c(n)-B_{r+1}^c(-n)}
  {2^{r+1}(r+1)r!}.
\end{equation}
The generating function is regular at the origin and has the compact
form
\begin{equation}\label{eq:C-generating}
  \boxed{\;
  \sum_{r=0}^{\infty}C_r(n)t^r
  =\sum_{i=1}^{n}e^{(i-m_n)t}
  =\frac{\sinh(nt/2)}{\sinh(t/2)}.\;}
\end{equation}
Symmetry now gives the correct parity:
\[
  C_{2j+1}(n)=0.
\]
The first nonzero coefficients are
\[
  C_0(n)=n,\qquad
  C_2(n)=\frac{n(n^2-1)}{24},\qquad
  C_4(n)=\frac{n(n^2-1)(3n^2-7)}{5760}.
\]

\begin{theorem}[Midpoint expansion]\label{thm:midpoint}
Let \(f\in C^{2p+2}[1,n]\).  Then
\begin{equation}\label{eq:midpoint}
  \sum_{i=1}^{n}f(i)
  =
  \sum_{j=0}^{p}f^{(2j)}(m_n)C_{2j}(n)
  +\mathcal T_{2p+2},
\end{equation}
where
\[
  |\mathcal T_{2p+2}|
  \le
  \left\|f^{(2p+2)}\right\|_{\infty}
  C_{2p+2}(n).
\]
\end{theorem}
\begin{proof}
Apply Taylor's theorem at \(m_n\) to every \(f(i)\), sum the finite
expansions, and use \eqref{eq:C-definition}.  All odd powers cancel by
symmetry.  Summing the Lagrange remainder bounds gives the stated
estimate.
\end{proof}

For \(f(i)=\log a_i\), Theorem~\ref{thm:midpoint} is the rigorous
central Bernoulli bridge in MVDC.  The \(C_r(n)\) encode the centred
indices, while the derivatives of the smooth interpolation \(f(x)\)
encode the particular factors.  Consequently the same \(C_r(n)\) may
be reused for different smooth products without identifying them with
the data-dependent moments of the values \(f(i)\).

\section{Euler--Maclaurin form and the factorial}

For \(g_c(x)=f(x)-c\), the Euler--Maclaurin formula on the indices
\(1,\ldots,n\) is
\begin{align}
  \sum_{i=1}^{n}g_c(i)
  ={}&
  \int_1^n g_c(x)\,dx
  +\frac{g_c(1)+g_c(n)}2 \notag\\
  &+\sum_{j=1}^{p}
  \frac{B_{2j}}{(2j)!}
  \left(g_c^{(2j-1)}(n)-g_c^{(2j-1)}(1)\right)
  +\mathcal E_p . \label{eq:EM}
\end{align}
The standard remainder \(\mathcal E_p\) is an integral involving the
periodic Bernoulli polynomial and \(g_c^{(2p)}\).  A centre can absorb
some integral or boundary contributions in a specific application, but
those cancellations must be checked for that application.

For the factorial, choose the analytic centre
\[
  c_n=\log n-1+\frac{\log(2\pi n)}{2n},\qquad
  H_n=e^{nc_n}=\sqrt{2\pi n}\left(\frac ne\right)^n .
\]
Euler--Maclaurin applied to \(f(x)=\log x\), with the limiting constant
fixed in the standard way, yields
\begin{equation}\label{eq:factorial}
  \log\frac{n!}{H_n}
  \sim
  \sum_{j=1}^{\infty}
  \frac{B_{2j}}{2j(2j-1)n^{2j-1}}.
\end{equation}
Thus
\[
  \log\frac{n!}{H_n}
  =
  \frac1{12n}-\frac1{360n^3}
  +\frac1{1260n^5}-\frac1{1680n^7}
  +O(n^{-9}).
\]
These are precisely the analytic coefficients implemented in
\texttt{mvdc\_factorial\_analytic.py}.  For this particular centre, the
MVDC organization and the classical log--Stirling expansion are the
same closed formula written as \(H_n\) times a small correction; this
observation does not compare other MVDC centres or fitted cascades with
Stirling expansions.

\begin{table}[ht]
\centering
\begin{tabular}{rc}
\toprule
\(n\) & relative error after the \(n^{-7}\) term \\
\midrule
10  & \(8.23\times10^{-13}\) \\
20  & \(1.63\times10^{-15}\) \\
50  & \(4.31\times10^{-19}\) \\
100 & \(8.42\times10^{-22}\) \\
\bottomrule
\end{tabular}
\caption{High-precision check of \eqref{eq:factorial} using 80-digit arithmetic.}
\end{table}

\section{Application to truncated Euler products}

Let
\[
  \zeta_N(s)=\prod_{p\le N}(1-p^{-s})^{-1},
  \qquad \sigma:=\Re s>1.
\]
For \(M>N\), define
\[
  f_s(p):=-\log(1-p^{-s}),\qquad
  \mu_{N,M}(s):=\frac1n\sum_{N<p\le M}f_s(p),
  \qquad n=\pi(M)-\pi(N).
\]
The finite block is exactly
\[
  B_{N,M}(s):=\sum_{N<p\le M}f_s(p)=n\mu_{N,M}(s).
\]
Therefore the exact tail decomposition is
\begin{equation}\label{eq:zeta-tail}
  \log\frac{\zeta(s)}{\zeta_N(s)}
  =n\mu_{N,M}(s)+R_M(s),
  \qquad
  R_M(s)=\sum_{p>M}-\log(1-p^{-s}).
\end{equation}
The omitted tail has the elementary rigorous bound
\begin{equation}\label{eq:tail-bound}
  |R_M(s)|
  \le
  \frac{M^{1-\sigma}}
  {(\sigma-1)(1-2^{-\sigma})}.
\end{equation}
Indeed,
\(-\log(1-p^{-\sigma})\le p^{-\sigma}/(1-2^{-\sigma})\),
and the prime sum is bounded by the corresponding integer integral.

For numerical acceleration one may model \(R_M(2)\) by
\[
  I_M:=\int_M^\infty\frac{x^{-2}}{\log x}\,dx.
\]
This is a prime-density model, not the rigorous bound
\eqref{eq:tail-bound}.  With \(M=10N\), direct high-precision
calculation gives:

\begin{table}[ht]
\centering
\begin{tabular}{rrrrr}
\toprule
\(N\) & \(n\) & exact tail \(\Delta\) & \(B_{N,M}+I_M\) & error \\
\midrule
1000  & 1061 & \(1.26989814756\times10^{-4}\)
              & \(1.27049267685\times10^{-4}\) & \(5.95\times10^{-8}\) \\
5000  & 4464 & \(2.10858622405\times10^{-5}\)
              & \(2.10884676662\times10^{-5}\) & \(2.61\times10^{-9}\) \\
10000 & 8363 & \(9.81570113247\times10^{-6}\)
              & \(9.81712884998\times10^{-6}\) & \(1.43\times10^{-9}\) \\
\bottomrule
\end{tabular}
\caption{Exact finite MVDC block plus the integral model for the remaining zeta tail.}
\end{table}

The numerical accuracy in this table is genuine, but it comes from the
exact finite block and the integral tail model.  In the earlier
experiment the additional logarithmic central-moment term at \(N=1000\)
was \(-1.56\times10^{-14}\), far below the displayed
\(5.95\times10^{-8}\) error.  Such moments may still be studied as
diagnostics, but they are not an additional term in
\eqref{eq:zeta-tail}.

\subsection{Extension to Dirichlet \texorpdfstring{\(L\)}{L}-functions}

For a primitive Dirichlet character \(\chi\) modulo \(q\), define
\[
  f_{\chi,s}(p):=-\log(1-\chi(p)p^{-s}),\qquad p\nmid q.
\]
Using a branch of the logarithm consistent with the absolutely
convergent Euler product, the same exact decomposition holds:
\begin{equation}\label{eq:L-tail}
  \log\frac{L(s,\chi)}{L_N(s,\chi)}
  =
  n\mu_{N,M}(\chi,s)+R_{M,\chi}(s),
\end{equation}
where \(n\) counts the primes \(N<p\le M\) with \(p\nmid q\).  Since
\(|\chi(p)|\le1\), the bound \eqref{eq:tail-bound} also bounds
\(|R_{M,\chi}(s)|\).

For the nontrivial real character modulo \(3\),
\[
  \chi_3(n)=
  \begin{cases}
    1,&n\equiv1\pmod3,\\
   -1,&n\equiv2\pmod3,\\
    0,&3\mid n,
  \end{cases}
\]
the reference value can be computed independently from
\[
  L(2,\chi_3)
  =\frac1{9}
   \left[\zeta\!\left(2,\frac13\right)
        -\zeta\!\left(2,\frac23\right)\right].
\]
For \(N=10^3\) and \(M=10^6\), 60-digit arithmetic gives
\[
\begin{array}{lr}
\text{exact missing tail} & -5.61728190475735\times10^{-8},\\
\text{exact finite block }(N,M] & -5.61013233973667\times10^{-8},\\
\text{omitted tail }p>M & -7.14956502068784\times10^{-11}.
\end{array}
\]
Thus the previously observed \(10^{-10}\)-level agreement is retained,
with its source identified correctly as the explicitly computed prime
block rather than an added central-moment series.

\section{Conclusion}

The corrected analytic structure of MVDC has four distinct components:
\par\smallskip
\begin{itemize}
  \item the exact centre decomposition \eqref{eq:exact-centre}, valid for
        every positive finite product;
  \item the convergent factor-moment expansion
        \eqref{eq:factor-moments}, with an explicit remainder;
  \item the finite \(C\)--Bernoulli coefficients
        \eqref{eq:C-definition}, which encode centred index sums and
        enter smooth midpoint expansions through
        Theorem~\ref{thm:midpoint};
  \item application-specific analytic or numerical models for the
        residual, such as \eqref{eq:factorial} and
        \eqref{eq:zeta-tail}.
\end{itemize}
Central Bernoulli and Nörlund polynomials therefore remain a genuine
part of the MVDC framework.  What must not be conflated are centred
index coefficients, moments of factor deviations, and central moments
of logarithmic data.  Once those objects are kept separate, the valid
factorial and Euler-product computations remain intact and their error
sources can be stated precisely.

\begin{thebibliography}{9}
\bibitem{DLMF-EM}
NIST Digital Library of Mathematical Functions,
\emph{Euler--Maclaurin Formula}, Sec.~2.10(i),
\url{https://dlmf.nist.gov/2.10}.

\bibitem{DLMF-Bernoulli}
NIST Digital Library of Mathematical Functions,
\emph{Bernoulli and Euler Polynomials: Sums of Powers}, Sec.~24.4(iii),
\url{https://dlmf.nist.gov/24.4}.

\bibitem{DLMF-Generalized}
NIST Digital Library of Mathematical Functions,
\emph{Generalized Bernoulli and Nörlund Polynomials}, Sec.~24.16,
\url{https://dlmf.nist.gov/24.16}.

\bibitem{Tempesta}
P.~Tempesta,
\emph{On Appell sequences of polynomials of Bernoulli and Euler type},
\emph{Journal of Mathematical Analysis and Applications}
\textbf{341} (2008), 1295--1310.

\bibitem{Norlund}
N.~E.~Nörlund,
\emph{Vorlesungen über Differenzenrechnung},
Springer, Berlin, 1924.
\end{thebibliography}

\end{document}
